% Ch8 WS2 -- bond energies and enthalpy from bonds. Block 3.
\documentclass{apchem-worksheet}

\apcunitword{Chapter}
\apcunit{8}
\apcunittitle{Bonding: General Concepts}
\apcdoctitle{Worksheet 2 \textbullet\ Bond Energies}
\apcsubtitle{Zumdahl \S8.8 \textbullet\ broken minus formed --- the reverse of the formation equation}

\begin{document}
\wsheader[Bond energies (kJ/mol) for all problems:
\ce{H-H} 432 \textbullet\ \ce{F-F} 154 \textbullet\ \ce{Cl-Cl} 239
\textbullet\ \ce{Br-Br} 193 \textbullet\ \ce{H-F} 565 \textbullet\
\ce{H-Cl} 427 \textbullet\ \ce{H-Br} 363 \textbullet\ \ce{N#N} 941
\textbullet\ \ce{N-H} 391 \textbullet\ \ce{O=O} 495 \textbullet\
\ce{O-H} 467 \textbullet\ \ce{C-H} 413 \textbullet\ \ce{C-C} 347
\textbullet\ \ce{C=C} 614 \textbullet\ \ce{C#C} 839.]

\problem[]{Bond order, length, and energy:}
\begin{parts}
  \item Rank \ce{C-C}, \ce{C=C}, \ce{C#C} by increasing bond length.
        \answerblank[5cm]{\ce{C#C} $<$ \ce{C=C} $<$ \ce{C-C}}
  \item Is a triple bond three times as strong as a single bond? Support
        with the data. \answerlines[2]{No --- $839$ versus $347$ is about
        2.4 times, not 3. The increments shrink as bond order rises.}
  \item Explain why breaking a bond is always endothermic.
        \answerlines[2]{The bonded arrangement is lower in energy than
        separated atoms, so energy must be supplied to climb out of that
        well.}
\end{parts}

\problem[]{Compute $\Delta H$ from bond energies:}
\begin{parts}
  \item \ce{H2(g) + Br2(g) -> 2 HBr(g)} \workspace[2cm]
        \keyonly{\begin{apcanswer}Broken $432 + 193 = 625$; formed
        $2(363) = 726$; $\Delta H = 625 - 726 =
        \SI{-101}{\kilo\joule}$\end{apcanswer}}
  \item \ce{N2(g) + 3 H2(g) -> 2 NH3(g)} \workspace[2.2cm]
        \keyonly{\begin{apcanswer}Broken $941 + 3(432) = 2237$; formed
        $6(391) = 2346$; $\Delta H =
        \SI{-109}{\kilo\joule}$\end{apcanswer}}
  \item \ce{2 H2(g) + O2(g) -> 2 H2O(g)} \workspace[2.2cm]
        \keyonly{\begin{apcanswer}Broken $2(432) + 495 = 1359$; formed
        $4(467) = 1868$; $\Delta H =
        \SI{-509}{\kilo\joule}$\end{apcanswer}}
\end{parts}

\problem[]{The accepted value for problem 2(b) is
\SI{-92.2}{\kilo\joule}, but bond energies give \SI{-109}{\kilo\joule}.}
\begin{parts}
  \item Explain the discrepancy. \answerlines[2]{Tabulated bond energies are
        averages taken across many different compounds, so any calculation
        built from them yields an estimate rather than an exact value.}
  \item Would you expect better or worse agreement for
        \ce{H2 + F2 -> 2 HF}, where every species is diatomic? Explain.
        \answerlines[2]{Better --- for diatomic molecules the bond energy is
        a directly measured value for that specific bond, not an average, so
        the estimate is far more reliable.}
\end{parts}

\problem[]{Compute $\Delta H$ for the hydrogenation of ethene,
\ce{C2H4(g) + H2(g) -> C2H6(g)}. Ethene contains one \ce{C=C} and four
\ce{C-H}; ethane contains one \ce{C-C} and six \ce{C-H}.
\workspace[3cm]
\keyonly{\begin{apcanswer}Broken: \ce{C=C} 614 $+$ 4 \ce{C-H} $1652$ $+$
\ce{H-H} 432 $= 2698$. Formed: \ce{C-C} 347 $+$ 6 \ce{C-H} $2478 = 2825$.
$\Delta H = 2698 - 2825 = \SI{-127}{\kilo\joule}$.

Shortcut: only the C=C$\to$C--C change and two new C--H bonds matter:
$(614 + 432) - (347 + 2 \times 413) = 1046 - 1173 =
\SI{-127}{\kilo\joule}$\end{apcanswer}}}

\problem[]{A student computes $\Delta H$ as ``bonds formed minus bonds
broken'' and reports \SI{+109}{\kilo\joule} for ammonia synthesis. Explain
the error and how to catch it without redoing the arithmetic.
\answerlines[3]{The correct order is broken minus formed --- reversing it
flips the sign. The sanity check: the products' bonds (2346 kJ) are
collectively stronger than the reactants' (2237 kJ), so energy must be
released and $\Delta H$ must be negative.}}

\problem[]{Predict, without calculating, whether
\ce{2 HCl(g) -> H2(g) + Cl2(g)} is exothermic or endothermic. Justify from
the bond energies alone. \answerlines[2]{Endothermic. Breaking two strong
H--Cl bonds ($2 \times 427 = 854$) costs more than forming one H--H and one
Cl--Cl bond releases ($432 + 239 = 671$).}}

\begin{spiralstrip}{Chapter 8 \textbullet\ blocks 1--2}
  \item Most polar bond: \ce{C-N}, \ce{C-O}, or \ce{C-F}?
        \answerblank[2.4cm]{\ce{C-F}}
  \item Larger: \ce{S^2-} or \ce{Cl-}? \answerblank[2.4cm]{\ce{S^2-}}
  \item Higher lattice energy: \ce{NaF} or \ce{MgO}?
        \answerblank[2.4cm]{\ce{MgO}}
  \item Is \ce{CCl4} polar? \answerblank[2cm]{no}
  \item Lattice energy depends on charges and:
        \answerblank[2.6cm]{ionic radii}
\end{spiralstrip}

\end{document}
